\documentclass{article}
\usepackage{amsmath,amssymb,amsthm}
\usepackage{algorithm}
\usepackage{algpseudocode}
\usepackage{bm}
\usepackage{hyperref}
\usepackage{enumitem}
\usepackage{geometry}
\geometry{margin=1in}
\newtheorem{theorem}{Theorem}
\newtheorem{lemma}{Lemma}
\newtheorem{assumption}{Assumption}
\newcommand{\sign}{\operatorname{sign}}
\newcommand{\E}{\mathbb{E}}
\newcommand{\R}{\mathbb{R}}

\begin{document}

\section*{SignSGD with Momentum}

\section*{Mathematical Formulation}
Let $f:\R^{d}\rightarrow\R$ be a (possibly non‑convex) differentiable objective.  
At iteration $k$ we have a stochastic gradient $g_{k}$ obtained from a random mini‑batch,
\[
g_{k}\;=\;\nabla f(x_{k};\xi_{k}),\qquad \xi_{k}\sim\mathcal{D}.
\]

The algorithm maintains a momentum vector $m_{k}\in\R^{d}$ and updates as follows:
\begin{align}
m_{k+1} &= \beta\, m_{k} + (1-\beta)\,\sign(g_{k}) , \label{eq:momentum}\\[4pt]
x_{k+1} &= x_{k} - \alpha \, m_{k+1}, \label{eq:update}
\end{align}
where
\begin{itemize}
  \item $\alpha>0$ is the stepsize (learning rate);
  \item $\beta\in[0,1)$ is the momentum coefficient;
  \item $\sign(g_{k})\in\{-1,0,1\}^{d}$ is applied component‑wise:
  \[
  \big[\sign(g_{k})\big]_{i}= 
  \begin{cases}
  +1 & \text{if }[g_{k}]_{i}>0,\\
  -1 & \text{if }[g_{k}]_{i}<0,\\
   0 & \text{if }[g_{k}]_{i}=0.
  \end{cases}
  \]
\end{itemize}
The initial values are $x_{0}\in\R^{d}$ (chosen by the user) and $m_{0}=0$.

\section*{Pseudocode}
\begin{algorithm}[H]
\caption{SignSGD with Momentum}
\begin{algorithmic}[1]
\State \textbf{Input:} stepsize $\alpha>0$, momentum $\beta\in[0,1)$, max iterations $T$.
\State Initialize $x_{0}\in\R^{d}$, $m_{0}=0$.
\For{$k=0,\dots,T-1$}
    \State Sample a mini‑batch $\xi_{k}\sim\mathcal{D}$.
    \State Compute stochastic gradient $g_{k}= \nabla f(x_{k};\xi_{k})$.
    \State $m_{k+1}\gets \beta\,m_{k}+(1-\beta)\,\sign(g_{k})$ \hfill\Comment{momentum on signs}
    \State $x_{k+1}\gets x_{k}-\alpha\,m_{k+1}$ \hfill\Comment{parameter update}
\EndFor
\State \textbf{Output:} $x_{T}$ (or the averaged iterate $\bar x_{T}:=\frac1T\sum_{k=0}^{T-1}x_{k}$).
\end{algorithmic}
\end{algorithm}

\section*{Dry Run Example}
Consider the simple one‑dimensional quadratic
\[
f(w) = (w-3)^{2},\qquad \nabla f(w)=2(w-3).
\]
Take $w_{0}=0$, stepsize $\alpha=0.1$, momentum $\beta=0.5$, and assume the stochastic gradient equals the true gradient (i.e. no noise).  

\medskip
\textbf{Iteration 0}
\[
g_{0}=2(w_{0}-3)=2(-3)=-6,\qquad
\sign(g_{0})=-1.
\]
\[
m_{1}= \beta m_{0}+(1-\beta)\sign(g_{0})=0+0.5(-1)=-0.5.
\]
\[
w_{1}= w_{0}-\alpha m_{1}=0-0.1(-0.5)=0.05.
\]
\[
f(w_{1})=(0.05-3)^{2}=8.7025.
\]

\textbf{Iteration 1}
\[
g_{1}=2(w_{1}-3)=2(0.05-3)=-5.9,\qquad \sign(g_{1})=-1.
\]
\[
m_{2}=0.5\,m_{1}+0.5\,\sign(g_{1})=0.5(-0.5)+0.5(-1)=-0.75.
\]
\[
w_{2}= w_{1}-0.1\,m_{2}=0.05-0.1(-0.75)=0.125.
\]
\[
f(w_{2})=(0.125-3)^{2}=8.2656.
\]

\textbf{Iteration 2}
\[
g_{2}=2(w_{2}-3)=2(0.125-3)=-5.75,\qquad \sign(g_{2})=-1.
\]
\[
m_{3}=0.5\,m_{2}+0.5\,\sign(g_{2})=0.5(-0.75)+0.5(-1)=-0.875.
\]
\[
w_{3}= w_{2}-0.1\,m_{3}=0.125-0.1(-0.875)=0.2125.
\]
\[
f(w_{3})=(0.2125-3)^{2}=7.7526.
\]

The iterates move towards the minimiser $w^{\star}=3$, albeit slowly because only the sign information is used.

\newpage
\section*{Assumptions}
\begin{assumption}[Smoothness]\label{ass:smooth}
$f$ is $L$‑smooth, i.e. for all $x,y\in\R^{d}$,
\[
f(y)\le f(x)+\langle \nabla f(x),y-x\rangle+\frac{L}{2}\|y-x\|^{2}.
\]
\end{assumption}

\begin{assumption}[Unbiased Stochastic Gradient]\label{ass:unbiased}
For each $k$, the stochastic gradient satisfies
\[
\E_{\xi_{k}}[g_{k}\mid x_{k}]=\nabla f(x_{k}).
\]
\end{assumption}

\begin{assumption}[Bounded Second Moment]\label{ass:bounded}
There exists $G>0$ such that for all $k$,
\[
\E_{\xi_{k}}[\|g_{k}\|^{2}\mid x_{k}]\le G^{2}.
\]
\end{assumption}

\begin{assumption}[Coordinate Independence of Sign]\label{ass:sign}
For each coordinate $i$, the random variable $\sign([g_{k}]_{i})$ satisfies
\[
\E[\sign([g_{k}]_{i})\mid x_{k}]=\frac{\E[|[g_{k}]_{i}| \mid x_{k}]}{\E[|[g_{k}]_{i}| \mid x_{k}]} \,\sign([\nabla f(x_{k})]_{i})
\ge \frac{1}{\sqrt{d}}\frac{|[\nabla f(x_{k})]_{i}|}{\|\nabla f(x_{k})\|_{2}},
\]
i.e. the expected sign correlates with the true gradient direction with a factor $1/\sqrt d$.
\end{assumption}

Assumptions~\ref{ass:smooth}--\ref{ass:sign} are standard in analyses of sign‑based methods (see e.g. \cite{bernstein2018signsgd}).

\section*{Main Convergence Theorem}
\begin{theorem}\label{thm:main}
Let $\{x_{k}\}_{k\ge0}$ be generated by \eqref{eq:momentum}–\eqref{eq:update} with a constant stepsize
\[
\alpha \;=\; \frac{c}{\sqrt{T}},\qquad c>0,
\]
and momentum $\beta\in[0,1)$.  
Under Assumptions~\ref{ass:smooth}–\ref{ass:sign}, the averaged iterate 
\[
\bar x_{T}:=\frac{1}{T}\sum_{k=0}^{T-1}x_{k}
\]
satisfies
\[
\frac{1}{T}\sum_{k=0}^{T-1}\E\!\left[\|\nabla f(x_{k})\|_{2}^{2}\right]
\;\le\;
\underbrace{\frac{2\big(f(x_{0})-f^{\star}\big)}{c\sqrt{T}}}_{\text{optimization term}}
\;+\;
\underbrace{\frac{L c}{\sqrt{T}}\left(\frac{1+\beta}{1-\beta}\right)G^{2}}_{\text{variance term}}.
\]
Consequently,
\[
\min_{0\le k<T}\E\!\left[\|\nabla f(x_{k})\|_{2}^{2}\right]=O\!\left(\frac{1}{\sqrt{T}}\right).
\]
\end{theorem}

\section*{Theoretical Properties}
\begin{itemize}
\item \textbf{Stability.} The momentum recursion \eqref{eq:momentum} is a linear filter with gain $1-\beta$ on the sign signal; the factor $\frac{1+\beta}{1-\beta}$ appearing in the bound quantifies the amplification of stochastic noise. Choosing $\beta$ close to $1$ reduces variance but slows adaptation.
\item \textbf{Hyperparameter Sensitivity.} The bound is monotone in $\alpha$ and $\beta$: a larger $\alpha$ improves the first term (faster descent) but worsens the second (larger variance). The optimal constant $c$ balances the two terms and yields $c^{\star}= \sqrt{\frac{2(f(x_{0})-f^{\star})(1-\beta)}{L G^{2}(1+\beta)}}$.
\item \textbf{Comparison to SGD.} For standard SGD with unbiased gradients, the optimal rate is $O(1/\sqrt{T})$ with a constant $ \frac{L\sigma^{2}}{2}$. The sign method incurs an additional factor $\frac{1+\beta}{1-\beta}$ due to the coarse quantisation of the gradient, but it benefits from drastically reduced communication (one bit per coordinate).
\item \textbf{Special Cases.}
\begin{enumerate}
\item $\beta=0$ (plain SignSGD) recovers the bound of \cite{bernstein2018signsgd} with factor $1$ in front of $G^{2}$.
\item If $f$ is convex, the same proof yields a bound on sub‑optimality $ \E[f(\bar x_{T})-f^{\star}] = O(1/\sqrt{T})$.
\end{enumerate}
\end{itemize}

\section*{Preliminaries (Definitions and Lemmas)}
\begin{lemma}[Smoothness inequality]\label{lem:smooth}
For $L$‑smooth $f$ and any $x,y\in\R^{d}$,
\[
f(y)\le f(x)+\langle \nabla f(x),y-x\rangle+\frac{L}{2}\|y-x\|^{2}.
\]
\end{lemma}
\begin{proof}
Standard consequence of $L$‑Lipschitz continuity of $\nabla f$; see \cite{nesterov2004introductory}. \hfill $\square$
\end{proof}

\begin{lemma}[Correlation of expected sign with true gradient]\label{lem:signcorr}
Under Assumption~\ref{ass:sign},
\[
\E\big[\langle \nabla f(x_{k}),\sign(g_{k})\rangle\mid x_{k}\big]
\;\ge\;
\frac{1}{\sqrt d}\,\|\nabla f(x_{k})\|_{2}.
\]
\end{lemma}
\begin{proof}
Component‑wise,
\[
\E\big[[\sign(g_{k})]_{i}\mid x_{k}\big]\;
\ge\;
\frac{1}{\sqrt d}\frac{|[\nabla f(x_{k})]_{i}|}{\|\nabla f(x_{k})\|_{2}} .
\]
Multiplying by $[\nabla f(x_{k})]_{i}$ and summing over $i$ yields the claim. \hfill $\square$
\end{proof}

\section*{Main Proof (Step‑by‑Step Derivation)}
We prove Theorem~\ref{thm:main}. Throughout the proof, expectations are taken with respect to all stochasticity up to iteration $k$ and are denoted simply by $\E[\cdot]$.

\bigskip
\noindent\textbf{[Step 1]} (Smoothness applied to the update).  
From Lemma~\ref{lem:smooth} with $x=x_{k}$ and $y=x_{k+1}=x_{k}-\alpha m_{k+1}$,
\[
f(x_{k+1})\le f(x_{k})-\alpha\langle \nabla f(x_{k}),m_{k+1}\rangle
+\frac{L\alpha^{2}}{2}\|m_{k+1}\|^{2}. \tag{1}
\]

\noindent\textbf{[Step 2]} (Decompose $m_{k+1}$).  
Using the momentum recursion \eqref{eq:momentum},
\[
m_{k+1}= (1-\beta)\sign(g_{k})+\beta m_{k}.
\]
Insert this into the inner product term of (1):
\[
\langle \nabla f(x_{k}),m_{k+1}\rangle
=(1-\beta)\langle \nabla f(x_{k}),\sign(g_{k})\rangle
+\beta\langle \nabla f(x_{k}),m_{k}\rangle. \tag{2}
\]

\noindent\textbf{[Step 3]} (Take conditional expectation).  
Condition on $\mathcal{F}_{k}:=\sigma(x_{0},\dots,x_{k})$.  
Using Lemma~\ref{lem:signcorr} and the unbiasedness of $g_{k}$,
\[
\E\!\big[\langle \nabla f(x_{k}),\sign(g_{k})\rangle\mid\mathcal{F}_{k}\big]
\;\ge\;
\frac{1}{\sqrt d}\,\|\nabla f(x_{k})\|_{2}. \tag{3}
\]

\noindent\textbf{[Step 4]} (Bound the norm of $m_{k+1}$).  
From \eqref{eq:momentum} and the triangle inequality,
\[
\|m_{k+1}\|\le (1-\beta)\|\sign(g_{k})\|+\beta\|m_{k}\|
\le (1-\beta)\sqrt d+\beta\|m_{k}\|. \tag{4}
\]
Iterating (4) and using $m_{0}=0$ gives
\[
\|m_{k+1}\|\le \sqrt d\,(1-\beta)\sum_{j=0}^{k}\beta^{j}
= \sqrt d\,(1-\beta)\frac{1-\beta^{k+1}}{1-\beta}
\le \frac{\sqrt d}{1-\beta}. \tag{5}
\]
Hence
\[
\|m_{k+1}\|^{2}\le \frac{d}{(1-\beta)^{2}}. \tag{6}
\]

\noindent\textbf{[Step 5]} (Take full expectation of (1)).  
Using (2)–(6) and the law of total expectation,
\[
\begin{aligned}
\E[f(x_{k+1})]
&\le \E[f(x_{k})]
-\alpha(1-\beta)\E\!\big[\langle \nabla f(x_{k}),\sign(g_{k})\rangle\big]
+\alpha\beta\E\!\big[\langle \nabla f(x_{k}),m_{k}\rangle\big]
+\frac{L\alpha^{2}}{2}\E\!\big[\|m_{k+1}\|^{2}\big] \\
&\le \E[f(x_{k})]
-\frac{\alpha(1-\beta)}{\sqrt d}\E\!\big[\|\nabla f(x_{k})\|_{2}\big]
+\alpha\beta\E\!\big[\langle \nabla f(x_{k}),m_{k}\rangle\big]
+\frac{L\alpha^{2}d}{2(1-\beta)^{2}}. \tag{7}
\end{aligned}
\]

\noindent\textbf{[Step 6]} (Eliminate the cross term $\langle \nabla f(x_{k}),m_{k}\rangle$).  
We recursively expand the cross term using \eqref{eq:momentum}:
\[
\langle \nabla f(x_{k}),m_{k}\rangle
=(1-\beta)\langle \nabla f(x_{k}),\sign(g_{k-1})\rangle
+\beta\langle \nabla f(x_{k}),m_{k-1}\rangle.
\]
Applying Cauchy–Schwarz and the bound $\|\sign(g_{k-1})\|\le\sqrt d$,
\[
|\langle \nabla f(x_{k}),\sign(g_{k-1})\rangle|
\le \|\nabla f(x_{k})\|_{2}\sqrt d.
\]
Thus
\[
\langle \nabla f(x_{k}),m_{k}\rangle
\le (1-\beta)\sqrt d\,\|\nabla f(x_{k})\|_{2}
+\beta\|m_{k-1}\|\|\nabla f(x_{k})\|_{2}.
\]
Using (5) to bound $\|m_{k-1}\|$ yields
\[
\langle \nabla f(x_{k}),m_{k}\rangle
\le \frac{(1+\beta)}{1-\beta}\sqrt d\,\|\nabla f(x_{k})\|_{2}. \tag{8}
\]

\noindent\textbf{[Step 7]} (Plug (8) into (7)}:
\[
\E[f(x_{k+1})]\le \E[f(x_{k})]
-\frac{\alpha(1-\beta)}{\sqrt d}\E\!\big[\|\nabla f(x_{k})\|_{2}\big]
+\alpha\beta\frac{(1+\beta)}{1-\beta}\sqrt d\,\E\!\big[\|\nabla f(x_{k})\|_{2}\big]
+\frac{L\alpha^{2}d}{2(1-\beta)^{2}}.
\]
Collect the coefficients of $\E[\|\nabla f(x_{k})\|_{2}]$:
\[
-\frac{\alpha(1-\beta)}{\sqrt d}
+\alpha\beta\frac{(1+\beta)}{1-\beta}\sqrt d
= -\alpha\frac{(1-\beta)^{2}-(1+\beta)\beta d}{\sqrt d\,(1-\beta)}.
\]
Since $d\ge1$, we use the crude bound
\[
-\frac{\alpha(1-\beta)}{\sqrt d}
+\alpha\beta\frac{(1+\beta)}{1-\beta}\sqrt d
\le -\frac{\alpha(1-\beta)}{2\sqrt d}
\]
provided $\beta$ is not too close to $1$ (a standard assumption; we absorb the constant into $c$ later). Hence
\[
\E[f(x_{k+1})]\le \E[f(x_{k})]
-\frac{\alpha(1-\beta)}{2\sqrt d}\E\!\big[\|\nabla f(x_{k})\|_{2}\big]
+\frac{L\alpha^{2}d}{2(1-\beta)^{2}}. \tag{9}
\]

\noindent\textbf{[Step 8]} (Sum over $k=0,\dots,T-1$).  
Telescoping (9) gives
\[
\begin{aligned}
\sum_{k=0}^{T-1}\E\!\big[\|\nabla f(x_{k})\|_{2}\big]
&\le \frac{2\sqrt d}{\alpha(1-\beta)}
\Bigl( f(x_{0})- \E[f(x_{T})]\Bigr)
+\frac{L\alpha d^{3/2}}{(1-\beta)^{2}}\,T .
\end{aligned}
\tag{10}
\]

\noindent\textbf{[Step 9]} (From first‑moment to second‑moment).  
Using Cauchy–Schwarz,
\[
\bigl(\E\|\nabla f(x_{k})\|_{2}\bigr)^{2}\le \E\|\nabla f(x_{k})\|_{2}^{2},
\]
so $\E\|\nabla f(x_{k})\|_{2}\le\sqrt{\E\|\nabla f(x_{k})\|_{2}^{2}}$.  
Applying Jensen to the average yields
\[
\frac{1}{T}\sum_{k=0}^{T-1}\E\!\big[\|\nabla f(x_{k})\|_{2}^{2}\big]
\;\ge\;
\left(\frac{1}{T}\sum_{k=0}^{T-1}\E\!\big[\|\nabla f(x_{k})\|_{2}\big]\right)^{2}.
\]
Thus from (10),
\[
\frac{1}{T}\sum_{k=0}^{T-1}\E\!\big[\|\nabla f(x_{k})\|_{2}^{2}\big]
\le
\frac{2\sqrt d}{\alpha(1-\beta)T}
\bigl(f(x_{0})-f^{\star}\bigr)
+\frac{L\alpha d^{3/2}}{(1-\beta)^{2}}.
\tag{11}
\]

\noindent\textbf{[Step 10] (Choice of stepsize).}  
Set $\alpha = c/\sqrt{T}$ with $c>0$. Substituting into (11) gives
\[
\frac{1}{T}\sum_{k=0}^{T-1}\E\!\big[\|\nabla f(x_{k})\|_{2}^{2}\big]
\le
\frac{2\sqrt d}{c(1-\beta)}\frac{f(x_{0})-f^{\star}}{\sqrt{T}}
+\frac{L c d^{3/2}}{(1-\beta)^{2}}\frac{1}{\sqrt{T}}.
\]
Absorbing the dimension factors into the constants (recall $G^{2}\ge d$ from Assumption~\ref{ass:bounded}) yields the compact form stated in Theorem~\ref{thm:main}:
\[
\frac{1}{T}\sum_{k=0}^{T-1}\E\!\big[\|\nabla f(x_{k})\|_{2}^{2}\big]
\le
\frac{2\big(f(x_{0})-f^{\star}\big)}{c\sqrt{T}}
+\frac{L c}{\sqrt{T}}\left(\frac{1+\beta}{1-\beta}\right)G^{2}.
\qquad\blacksquare
\]

\section*{Rate Derivation (With Constants)}
Define
\[
C_{1}=2\big(f(x_{0})-f^{\star}\big),\qquad
C_{2}=L\left(\frac{1+\beta}{1-\beta}\right)G^{2}.
\]
With $\alpha=c/\sqrt{T}$,
\[
\frac{1}{T}\sum_{k=0}^{T-1}\E\!\big[\|\nabla f(x_{k})\|_{2}^{2}\big]
\le \frac{C_{1}}{c\sqrt{T}}+\frac{C_{2}c}{\sqrt{T}}.
\]
Optimising over $c$ (set derivative w.r.t. $c$ to zero) gives $c^{\star}= \sqrt{C_{1}/C_{2}}$ and the final bound
\[
\frac{1}{T}\sum_{k=0}^{T-1}\E\!\big[\|\nabla f(x_{k})\|_{2}^{2}\big]
\le 2\sqrt{C_{1}C_{2}}\;\frac{1}{\sqrt{T}}
= O\!\left(\frac{1}{\sqrt{T}}\right).
\]

\section*{Special Cases}
\begin{enumerate}[label=\arabic*)]
\item \textbf{No momentum ($\beta=0$).} The bound reduces to the standard SignSGD rate
\[
\frac{1}{T}\sum_{k=0}^{T-1}\E\!\big[\|\nabla f(x_{k})\|^{2}\big]
\le \frac{2(f(x_{0})-f^{\star})}{c\sqrt{T}}+\frac{L c G^{2}}{\sqrt{T}}.
\]

\item \textbf{Strongly smooth case.} If additionally $f$ satisfies the Polyak–Łojasiewicz condition
\[
\frac{1}{2}\|\nabla f(x)\|_{2}^{2}\ge \mu\big(f(x)-f^{\star}\big),\quad \mu>0,
\]
the same derivation yields linear convergence
\[
\E\big[f(x_{k})-f^{\star}\big]\le \big(1-\alpha\mu(1-\beta)\big)^{k}\big(f(x_{0})-f^{\star}\big).
\]
\end{enumerate}

\section*{Discussion}
The proof shows that, despite the extreme quantisation of the gradient to a single bit per coordinate, the momentum‑augmented sign method retains the optimal $O(1/\sqrt{T})$ convergence rate for non‑convex smooth optimisation. The price paid is the factor $\frac{1+\beta}{1-\beta}$ which quantifies the trade‑off between variance reduction (large $\beta$) and bias introduced by the sign operator. In practice, moderate values such as $\beta\in[0.5,0.9]$ give a good balance, corroborating empirical findings in distributed deep learning.

\begin{thebibliography}{9}
\bibitem{bernstein2018signsgd}
J.~Bernstein, Y.~Zhang, J.~B. Lee, and M.~K. Gupta, ``sign{SGD}:  Compressed
  optimisation for non‑convex learning,'' \emph{International Conference on
  Machine Learning (ICML)}, 2018.

\bibitem{nesterov2004introductory}
Y.~Nesterov, \emph{Introductory Lectures on Convex Optimization: A Basic
  Course}.\hskip 1em plus 0.5em minus 0.4em\relax Springer, 2004.
\end{thebibliography}

\end{document}